\chapter{Coronary Artery Bypass Grafting (CABG)}

\section*{Introduction \& Clinical Indications}
Coronary Artery Bypass Grafting (CABG) is a critical surgical intervention aimed at restoring blood flow to the heart muscle in patients suffering from significant coronary artery disease (CAD). This procedure involves the creation of new pathways for blood flow by bypassing obstructed coronary arteries using grafts, which can be either arterial or venous in origin. The most commonly utilized conduits include the internal mammary artery (IMA) and the great saphenous vein (GSV). CABG is particularly indicated in patients with multi-vessel disease, left main coronary artery stenosis, or those who have not responded adequately to medical therapy or percutaneous coronary intervention (PCI). The significance of CABG lies in its ability to alleviate symptoms such as angina, improve cardiac function, and enhance overall survival and quality of life for patients with advanced coronary artery disease.

\begin{itemize}
  \item Heart Bypass Surgery
  \item Coronary Bypass Surgery
  \item Open Heart Surgery
  \item Myocardial Revascularization Surgery
  \item Aortocoronary Bypass
  \item CABG Surgery
\end{itemize}

The primary surgical principle of CABG is to bypass obstructed coronary arteries to restore adequate blood flow to ischemic myocardial regions. Atherosclerosis leads to the narrowing of coronary arteries, resulting in reduced oxygen supply to the heart muscle, which can cause angina and myocardial infarction. CABG addresses this by creating new conduits that allow blood to flow from the aorta or subclavian artery directly to the coronary arteries beyond the blockages.

The rationale for CABG extends beyond symptom relief; it aims to improve myocardial viability and enhance left ventricular function, ultimately prolonging survival in patients with severe coronary artery disease. The choice of graft material is crucial; arterial grafts, particularly the LIMA, demonstrate superior long-term patency compared to venous grafts. The technical goals include achieving complete revascularization of all significant stenoses in viable myocardial segments, maximizing clinical benefits, and providing a robust blood supply to the heart.

The evolution of CABG has been marked by several key milestones. In the early 20th century, indirect methods of myocardial revascularization were explored, such as the Vineberg procedure in the 1940s. The introduction of cardiopulmonary bypass by John Gibbon in the 1950s revolutionized cardiac surgery, allowing for direct coronary artery interventions.

In 1964, Dr. Vasilii Kolesov performed the first successful internal mammary artery anastomosis to a coronary artery, laying the groundwork for arterial grafting. However, it was Dr. René Favaloro's pioneering work at the Cleveland Clinic in 1967 that established CABG as a standard treatment for coronary artery disease with his successful saphenous vein graft CABG. Over the decades, techniques have evolved, including the adoption of off-pump CABG and minimally invasive approaches, continually refining outcomes and expanding patient eligibility.

CABG is indicated for patients with significant coronary artery disease who meet specific criteria, often determined through a multidisciplinary Heart Team discussion. Indications include:

\begin{itemize}
  \item Severe angina pectoris refractory to medical therapy or PCI.
  \item Significant left main coronary artery stenosis (>50\%).
  \item Multi-vessel coronary artery disease, especially with proximal LAD involvement.
  \item Reduced left ventricular ejection fraction (LVEF) with viable myocardium.
  \item Acute coronary syndromes when PCI is not feasible or has failed.
  \item Post-myocardial infarction complications, such as ventricular septal rupture.
  \item Restenosis after previous PCI, particularly with complex lesions.
\end{itemize}

Standardized diagnostic scores, such as the SYNTAX score, may guide the decision-making process.

CABG is often the primary revascularization strategy for patients with severe coronary artery disease, particularly those with left main disease or multi-vessel involvement. Guidelines from major cardiology societies recommend CABG as the first-line treatment in these scenarios, especially when the SYNTAX score indicates high anatomical complexity.

In contrast, PCI may be considered for single-vessel disease or less complex two-vessel disease without diabetes. CABG can also serve as a secondary or salvage procedure in cases of failed PCI or recurrent ischemia. The decision to proceed with CABG or PCI should be made collaboratively, considering the patient's overall clinical picture, including comorbidities and preferences.

\section*{Relevant Surgical Anatomy}
A comprehensive understanding of the cardiac anatomy is essential for the successful execution of CABG. The coronary arteries, which supply blood to the myocardium, originate from the ascending aorta. The left main coronary artery bifurcates into the left anterior descending (LAD) artery and the circumflex artery, while the right coronary artery (RCA) supplies the right ventricle and inferior wall of the left ventricle. 

Key anatomical landmarks include:

\begin{itemize}
  \item **Left Internal Mammary Artery (LIMA):** Originates from the left subclavian artery and is commonly used as a graft for the LAD due to its excellent long-term patency.
  \item **Great Saphenous Vein (GSV):** Harvested from the leg, it serves as a conduit for bypassing other coronary arteries.
  \item **Radial Artery:** An alternative arterial graft, harvested from the forearm, used when collateral circulation is adequate.
  \item **Calot's Triangle:** Important for identifying the anatomy during graft harvesting.
  \item **McBurney's Point:** While primarily associated with appendectomy, knowledge of this landmark aids in avoiding injury to surrounding structures during thoracic procedures.
\end{itemize}

The surgical approach requires careful dissection to avoid injury to the phrenic nerve and to ensure adequate exposure of the heart and great vessels.

\section*{Preoperative Workup \& Preparation}
A thorough preoperative assessment is essential for optimizing patient outcomes. This includes:

\begin{itemize}
  \item **Comprehensive Medical Evaluation:** Detailed history and physical examination, including cardiac function assessment and pulmonary function tests.

  \item **Laboratory Tests:** Routine blood tests, including complete blood count, electrolyte panel, renal and liver function tests, and coagulation profile.

  \item **Medication Management:** Adjustments to antiplatelet and anticoagulant medications, with careful consideration of bleeding risks.

  \item **NPO Status:** Patients are instructed to be NPO for at least 6-8 hours prior to surgery.

  \item **Informed Consent:** Detailed discussions regarding the procedure, risks, benefits, and alternatives.

  \item **Preoperative Education:** Instructions on postoperative expectations, pain management, and the importance of smoking cessation.

  \item **Antibiotic Prophylaxis:** Administration of intravenous antibiotics shortly before the incision.

  \item **DVT Prophylaxis:** Initiation of mechanical prophylaxis to prevent thromboembolic events.
\end{itemize}

Certain conditions may contraindicate CABG or increase surgical risks.

\textbf{Absolute Contraindications:}
\begin{itemize}
  \item End-stage coronary artery disease with no viable myocardium.
  \item Severe irreversible neurological impairment.
  \item Uncontrolled severe systemic infection.
  \item Malignancy with a very short life expectancy.
  \item Severe coagulopathy or bleeding disorders.
  \item Refusal of blood products in high-risk patients.
\end{itemize}

\textbf{Relative Contraindications \& Precautions:}
\begin{itemize}
  \item Advanced age and frailty.
  \item Severe chronic obstructive pulmonary disease (COPD).
  \item Significant renal dysfunction.
  \item Previous cardiac surgery (redo CABG).
  \item Severe peripheral vascular disease.
  \item Uncontrolled diabetes mellitus.
  \item Severe left ventricular dysfunction without viable myocardium.
  \item Active substance abuse or non-compliance.
\end{itemize}

\section*{Surgical Technique \& Procedure}
CABG is typically performed under general anesthesia, with the patient positioned supine. The procedure involves several key stages:

\begin{itemize}
  \item **Anesthesia and Monitoring:** General anesthesia is induced, and arterial and central venous catheters are placed for monitoring.
  
  \item **Incisions and Graft Harvesting:** A median sternotomy is performed to access the heart. The LIMA is dissected from the chest wall, and if necessary, the GSV is harvested from the leg.

  \item **Cardiopulmonary Bypass (CPB) Initiation:** The patient is placed on CPB, with cannulas inserted into the ascending aorta and right atrium. The aorta is cross-clamped, and cardioplegia is administered to protect the myocardium.

  \item **Coronary Anastomoses:** Distal anastomoses are performed first, connecting the grafts to the coronary arteries beyond the blockages. The LIMA is typically anastomosed to the LAD, while SVGs are connected to other target arteries.

  \item **Weaning from CPB and Decannulation:** After confirming the patency of anastomoses, the heart is de-aired, and the aortic cross-clamp is removed. The patient is weaned off CPB, and protamine is administered to reverse heparin.

  \item **Closure:** Chest tubes are placed, and the sternum is closed with stainless steel wires. The muscle, fascia, and skin layers are then closed in anatomical fashion.
\end{itemize}

In off-pump CABG (OPCAB), the heart remains beating, and specialized stabilizers are used to immobilize the target coronary artery, avoiding the need for CPB.

\section*{Complications \& Risks}
CABG carries inherent risks, which can be categorized as follows:

\begin{itemize}
  \item \textbf{General Surgical Complications:}
    \begin{itemize}
      \item Bleeding requiring reoperation or transfusions.
      \item Infection, including surgical site infections and pneumonia.
      \item Anesthesia-related complications.
      \item Deep vein thrombosis and pulmonary embolism.
    \end{itemize}
  
  \item \textbf{Cardiac Complications:}
    \begin{itemize}
      \item Myocardial infarction due to graft occlusion.
      \item Arrhythmias, particularly atrial fibrillation.
      \item Low cardiac output syndrome.
      \item Pericardial effusion or tamponade.
      \item Graft failure or occlusion.
    \end{itemize}

  \item \textbf{Neurological Complications:}
    \begin{itemize}
      \item Stroke, ranging from 1-5\%.
      \item Transient ischemic attack (TIA).
      \item Post-perfusion syndrome.
      \item Nerve injury.
    \end{itemize}

  \item \textbf{Pulmonary Complications:}
    \begin{itemize}
      \item Pneumonia and atelectasis.
      \item Respiratory failure requiring prolonged ventilation.
    \end{itemize}

  \item \textbf{Renal Complications:}
    \begin{itemize}
      \item Acute kidney injury, especially in patients with pre-existing renal dysfunction.
    \end{itemize}

  \item \textbf{Vascular/Wound Complications:}
    \begin{itemize}
      \item Sternal instability or non-union.
      \item Leg wound complications from graft harvesting.
      \item Radial artery harvest complications.
    \end{itemize}

  \item \textbf{Mortality:} The overall mortality rate for elective CABG is typically 1-3\%, but can be higher in emergent cases or patients with significant comorbidities.
\end{itemize}

\section*{Postoperative Care \& Recovery}
Postoperative recovery following CABG is structured and involves several phases. Immediately after surgery, patients are monitored in the cardiothoracic intensive care unit (CTICU). The first 24-48 hours are critical, with continuous assessment of vital signs, cardiac rhythm, and urine output. Patients are typically intubated and mechanically ventilated, with extubation occurring as soon as they are stable and able to breathe independently.

During the first week, patients are encouraged to participate in early mobilization, including sitting up and short walks, to prevent complications such as pneumonia and deep vein thrombosis. Pain management is prioritized, often utilizing intravenous analgesics and patient-controlled analgesia.

By the end of the first week, most patients are transferred to a regular hospital floor, focusing on increasing mobility and managing pain with oral medications. Discharge typically occurs within 5-7 days post-surgery, with long-term recovery guided by a cardiac rehabilitation program. Full recovery, including the resolution of fatigue and return to pre-surgical activity levels, can take 2-3 months.

During CABG, surgeons document several key findings, including the number and location of diseased coronary arteries, the degree of stenosis, and the quality of distal vessels for anastomosis. The types of grafts used are recorded, along with any intraoperative complications.

Pathological examination of harvested tissues, such as the saphenous vein or radial artery, is performed to confirm vessel integrity. While native coronary arteries are not routinely biopsied, the presence of severe atherosclerosis is visually confirmed. In cases of concurrent procedures, such as ventricular aneurysmectomy, histopathological analysis is conducted to assess myocardial damage.

A successful postoperative course following CABG is characterized by a gradual improvement in cardiac function and quality of life. Patients typically experience stable hemodynamics and a decrease in pain levels. Early mobilization is encouraged, and patients should participate in breathing exercises to prevent pulmonary complications.

As recovery progresses, patients are expected to demonstrate improved exercise tolerance and resolution of angina symptoms. Surgical incisions should heal without signs of infection, and fatigue is expected to gradually improve over weeks to months. A normal postoperative course signifies effective revascularization and a reduced risk of future cardiac events.

Postoperative care and follow-up are essential for ensuring optimal recovery and long-term outcomes. Key components include:

\begin{itemize}
  \item **Early Surgical Follow-up:** Initial visits within 1-2 weeks post-discharge to assess wound healing and address concerns.

  \item **Cardiology Follow-up:** Regular visits starting around 4-6 weeks post-op for ongoing management and medication adjustments.

  \item **Cardiac Rehabilitation:** Enrollment in a structured program for supervised exercise and education on heart-healthy living.

  \item **Medication Adherence:** Lifelong antiplatelet therapy and management of comorbidities.

  \item **Lifestyle Modifications:** Encouragement of heart-healthy lifestyle changes, including diet and exercise.

  \item **Activity Resumption:** Gradual increase in physical activity, with full return to normal activities typically by 2-3 months.

  \item **Warning Signs:** Education on symptoms that warrant immediate medical attention, such as chest pain or signs of infection.
  
  \item **Imaging/Labs:** Routine monitoring of risk factors through blood tests and potential imaging studies if symptoms recur.
\end{itemize}

\section*{Outcomes \& Prognosis}
Coronary artery disease remains a leading cause of morbidity and mortality globally, with its prevalence increasing with age and risk factors such as hypertension, diabetes, and smoking. In the United States, approximately 18.2 million adults have CAD, and over 200,000 CABG procedures are performed annually. The typical CABG patient is male, over 60 years old, often with multiple comorbidities. Multi-vessel disease accounts for over 70\% of cases, with in-hospital mortality rates for elective CABG ranging from 1-3\%, but higher in emergent cases.

The prognosis following CABG is generally favorable, particularly for well-selected patients. Success rates for symptom relief are high, with over 80-90\% experiencing significant reduction in angina. Long-term survival rates are superior to medical therapy alone, with five-year survival rates ranging from 85-90\% and ten-year rates from 70-80\%. 

Graft patency is a critical factor, with arterial grafts showing superior long-term outcomes compared to saphenous vein grafts. Recurrence of angina or need for re-intervention can occur, but CABG provides durable revascularization. Prognostic factors influencing outcomes include age, left ventricular function, and completeness of revascularization, with landmark trials consistently demonstrating the long-term benefits of CABG in specific patient populations.

\section*{References}
\begin{enumerate}
  \item MedlinePlus. Coronary Artery Bypass Grafting (CABG). U.S. National Library of Medicine; 2024. Available from: \url{https://medlineplus.gov/coronaryarterybypass.html}
  
  \item Sabiston Textbook of Surgery: The Biological Basis of Modern Surgical Practice. 21st ed. Philadelphia, PA: Elsevier; 2021.
  
  \item American College of Cardiology/American Heart Association Task Force on Clinical Practice Guidelines. 2021 ACC/AHA Guideline for Coronary Artery Revascularization. Circulation. 2022;145(3):e4-e17.
  
  \item Farkouh ME, et al. Strategies for Multivessel Revascularization in Patients with Diabetes. N Engl J Med. 2012;367(25):2375-2384. (FREEDOM Trial)
\end{enumerate}

\clearpage
